\subsection{Теоретична рамка: інтеграція TPACK, студійної педагогіки та компетентнісного підходу}\label{subsec:1.4}

Теоретичне обґрунтування формування інформаційно-комунікативної компетентності (ІКК) майбутніх дизайнерів потребує інтегративної рамки, що поєднує три взаємопов'язані теоретичні перспективи: модель технологічних педагогічних та змістових знань (TPACK), студійну педагогіку та компетентнісний підхід. Кожна з цих перспектив висвітлює окремий аспект проблеми формування ІКК, а їх синтез забезпечує цілісне бачення, необхідне для аналізу даних у розділі~\ref{sec:2} та розробки моделі трансформації у розділі~\ref{sec:3}.

\subsubsection{Модель TPACK у контексті дизайн-освіти}\label{sssec:1.4.1}

Модель технологічних педагогічних та змістових знань (Technological Pedagogical Content Knowledge~--- TPACK), запропонована \citet{Mishra2006} на основі концепції педагогічних знань змісту \citet{Shulman1986}, описує взаємодію трьох базових компонентів знань викладача: змістових (Content Knowledge~--- CK), педагогічних (Pedagogical Knowledge~--- PK) та технологічних (Technological Knowledge~--- TK). На перетинах цих компонентів виникають інтегровані форми знань: технологічних змістових (TCK), технологічних педагогічних (TPK), педагогічних змістових (PCK) та центральний конструкт~--- TPACK, що поєднує всі три виміри.

В контексті дизайн-освіти ери штучного інтелекту модель TPACK набуває специфічного змісту:

\begin{itemize}
    \item \textbf{Змістові знання (CK)} охоплюють предметну область дизайну: теорію кольору, композицію, типографіку, історію мистецтва та дизайну, ергономіку, матеріалознавство. Це ядро фахової підготовки, що визначає \textit{що} дизайнер повинен знати;
    \item \textbf{Педагогічні знання (PK)} включають студійні методи навчання: майстер-класи, критики (crit sessions), проєктне навчання, рефлективну практику за \citet{Schon1983}. Це визначає \textit{як} відбувається навчання дизайну;
    \item \textbf{Технологічні знання (TK)} включають як традиційні цифрові інструменти дизайнера (Adobe Creative Suite, Figma, Blender, 3ds Max, AutoCAD), так і інструменти ШІ нового покоління (Midjourney, DALL-E, Stable Diffusion, Adobe Firefly, RunwayML) та навички промпт-інженерії. Це визначає \textit{якими засобами} працює сучасний дизайнер.
\end{itemize}

Ключова проблема, виявлена в оглядовому дослідженні~\cite{Musiienko2026ScopingReview}, полягає у системному дефіциті інтегрованих форм знань: лише 26,3\% досліджень генеративного ШІ у дизайн-освіті використовували будь-які теоретичні рамки, а домінування TPACK серед них свідчить про потребу саме в такому інтегративному підході. При цьому 73,7\% досліджень залишаються атеоретичними, що перешкоджає системному розумінню того, як інтеграція ШІ-інструментів впливає на педагогічні практики та зміст навчання.

Адаптація TPACK для формування ІКК дизайнерів передбачає переосмислення TK-компоненту: технологічні знання більше не обмежуються умінням працювати з конкретним програмним забезпеченням. Вони включають:
\begin{enumerate}
    \item здатність оцінювати та обирати оптимальний інструмент (традиційний чи ШІ-базований) для конкретного дизайнерського завдання;
    \item навички промпт-інженерії~--- формулювання текстових запитів для генеративних моделей, що забезпечують бажаний результат;
    \item критичне оцінювання результатів, згенерованих ШІ, з урахуванням естетичних, етичних та функціональних вимог;
    \item розуміння принципів роботи генеративних моделей на рівні, достатньому для ефективної взаємодії з ними.
\end{enumerate}

\subsubsection{Студійна педагогіка та рефлексивна практика}\label{sssec:1.4.2}

Студійна педагогіка (studio pedagogy) є <<підписовою педагогікою>> (signature pedagogy) дизайнерських професій, як визначив \citet{Shulman2005}. Її характерними рисами є:

\begin{itemize}
    \item \textit{навчання через діяльність} (learning by doing)~--- студенти набувають компетентностей через виконання дизайн-проєктів, а не через засвоєння теоретичних знань;
    \item \textit{критика (crit)}~--- публічне обговорення студентських робіт, що розвиває здатність до рефлексії, аргументації та комунікації;
    \item \textit{ітеративний процес}~--- дизайн передбачає багаторазове повторення циклу <<ідея~$\to$~прототип~$\to$~критика~$\to$~вдосконалення>>;
    \item \textit{формування професійної ідентичності}~--- студія є простором соціалізації, де формуються цінності та ставлення, характерні для професійної спільноти.
\end{itemize}

Теорія рефлексивної практики \citet{Schon1983} забезпечує теоретичне обґрунтування студійної педагогіки. Шон розрізняє два типи рефлексії: \textit{рефлексію-в-дії} (reflection-in-action), коли дизайнер аналізує та коригує свої дії в процесі роботи, та \textit{рефлексію-про-дію} (reflection-on-action), коли аналіз відбувається після завершення роботи. Обидва типи рефлексії є критично важливими для формування рефлексивного компоненту ІКК.

Генеративний ШІ трансформує студійну педагогіку в декількох напрямках. По-перше, змінюється роль студента: від виробника (maker), що створює артефакти своїми руками, до куратора (curator), що оцінює, обирає та редагує згенеровані ШІ варіанти~\cite{Musiienko2026ScopingReview}. По-друге, ітеративний процес значно прискорюється: ШІ-інструменти дозволяють генерувати десятки варіантів за хвилини, тоді як ручне створення потребувало б годин. По-третє, критика (crit) набуває нового виміру: студенти повинні обґрунтовувати не лише дизайнерські рішення, а й вибір інструментів, стратегію промптингу та критерії відбору серед згенерованих варіантів.

Ця трансформація не усуває потребу в ручних навичках (рисунок, живопис, макетування), а доповнює їх новими компетентностями, пов'язаними з ІКК. Студійна педагогіка, адаптована для ери ШІ, передбачає інтеграцію обох типів навичок у єдиний проєктний процес.

\subsubsection{Компетентнісний підхід та стандарт В2 Дизайн}\label{sssec:1.4.3}

Компетентнісний підхід, що став нормативним для європейського простору вищої освіти після Болонського процесу, визначає результати навчання через категорії знань, умінь та компетентностей~\cite{EuropeanCommission2017}. Європейська рамка кваліфікацій (EQF) встановлює тріадну категоризацію, що стала стандартом для розробки навчальних програм.

Стандарт вищої освіти України за спеціальністю 022 Дизайн рівня бакалавра (В2 Дизайн) визначає дві категорії компетентностей: загальні (ЗК) та фахові (ФК)~\cite{MON2018}. Аналіз, проведений у~\cite{Musiienko2025DesignEducation}, виявив структурну відповідність стандарту В2 Європейській рамці кваліфікацій за вимірами знань, умінь та автономності, проте з суттєвими прогалинами у компетентностях, пов'язаних із:
\begin{itemize}
    \item критичною рефлексією в контексті дизайнерської діяльності;
    \item міждисциплінарною співпрацею;
    \item сталим дизайном;
    \item цифровими та ШІ-навичками~--- що безпосередньо стосується ІКК.
\end{itemize}

Компетентності стандарту В2 можуть бути безпосередньо відображені на 4-компонентну структуру ІКК:

\begin{longtable}{@{}p{0.18\textwidth}p{0.35\textwidth}p{0.40\textwidth}@{}}
\caption{Відповідність компетентностей стандарту В2 Дизайн компонентам ІКК}\label{tab:b2-ikk-mapping} \\
\hline
\textbf{Компонент ІКК} & \textbf{Компетентності В2} & \textbf{Прогалина щодо ШІ} \\
\hline
\endfirsthead
\hline
\textbf{Компонент ІКК} & \textbf{Компетентності В2} & \textbf{Прогалина щодо ШІ} \\
\hline
\endhead
Інформаційно-аналітичний & ЗК2: здатність до пошуку та аналізу інформації з різних джерел & Відсутність згадки ШІ-генерованого контенту як джерела; не передбачено критичну оцінку ШІ-виходів \\
\addlinespace
Комунікативний & ЗК6: здатність до комунікації; ФК5: здатність до візуальної презентації & Не враховує цифрові комунікаційні платформи та AI-інструменти для презентації \\
\addlinespace
Технологічний & ФК8: здатність застосовувати комп'ютерні технології; ФК12: здатність використовувати сучасне ПЗ & Не згадує ШІ-інструменти, промпт-інженерію, генеративний дизайн \\
\addlinespace
Рефлексивний & ЗК1: здатність до критичного мислення; ЗК10: здатність до самоосвіти & Не передбачає рефлексію щодо етичних аспектів використання ШІ в дизайні \\
\hline
\end{longtable}

Таблиця~\ref{tab:b2-ikk-mapping} наочно демонструє, що стандарт В2 містить базові передумови для формування ІКК, проте не враховує ШІ-вимір, що є критичним у сучасних умовах. Це підтверджується даними аналізу 122 українських програм дизайну (п.~\ref{subsec:2.1}), де виявлено 0\% покриття компонентів <<Основи ШІ>>, <<Генеративний ШІ>> та <<Комп'ютерний зір>>.

\subsubsection{Інтегративна теоретична рамка}\label{sssec:1.4.4}

Синтез трьох розглянутих теоретичних перспектив утворює інтегративну рамку для дослідження формування ІКК майбутніх дизайнерів (рис.~\ref{fig:theoretical-framework}).

\begin{figure}[htbp]
\centering
\begin{tikzpicture}[
    node distance=2cm,
]
% Central node
\node[dia/header=12cm] (center) {ІНТЕГРАТИВНА ТЕОРЕТИЧНА РАМКА\\ФОРМУВАННЯ ІКК ДИЗАЙНЕРІВ};

% Three pillars
\node[dia/box=4.5cm, below left=2cm and 2cm of center, fill=blue!8] (tpack) {\textbf{TPACK}\\для дизайну};
\node[dia/box=4.5cm, below=2cm of center, fill=green!8] (studio) {\textbf{Студійна}\\педагогіка};
\node[dia/box=4.5cm, below right=2cm and 2cm of center, fill=orange!8] (competency) {\textbf{Компетентнісний}\\підхід (В2)};

% IKK components
\node[dia/rbox=3.5cm, below=2cm of tpack] (tech) {ТK $\to$ Технологічний\\компонент ІКК};
\node[dia/rbox=3.5cm, below=2cm of studio] (refl) {Рефлексія $\to$\\Рефлексивний\\компонент ІКК};
\node[dia/rbox=3.5cm, below=2cm of competency] (info) {ЗК/ФК $\to$ Інформаційно-\\аналітичний та\\Комунікативний\\компоненти ІКК};

% Arrows from center to pillars
\draw[dia/arrow] (center.south) -- (tpack.north);
\draw[dia/arrow] (center.south) -- (studio.north);
\draw[dia/arrow] (center.south) -- (competency.north);

% Arrows from pillars to IKK
\draw[dia/arrow] (tpack.south) -- (tech.north);
\draw[dia/arrow] (studio.south) -- (refl.north);
\draw[dia/arrow] (competency.south) -- (info.north);

% Cross-connections (dashed)
\draw[dia/line, dashed, gray] (tpack.east) -- (studio.west);
\draw[dia/line, dashed, gray] (studio.east) -- (competency.west);

% Bottom synthesis
\node[dia/header=12cm, below=1.5cm of refl] (bottom) {МОДЕЛЬ ТРАНСФОРМАЦІЇ\\НАВЧАЛЬНИХ ПРОГРАМ (Розділ~3)};
\draw[dia/arrow] (tech.south) -- (bottom.north);
\draw[dia/arrow] (refl.south) -- (bottom.north);
\draw[dia/arrow] (info.south) -- (bottom.north);

\end{tikzpicture}
\caption{Інтегративна теоретична рамка формування ІКК дизайнерів: синтез TPACK, студійної педагогіки та компетентнісного підходу}
\label{fig:theoretical-framework}
\end{figure}

Інтегративна рамка функціонує таким чином:

\begin{enumerate}
    \item \textbf{TPACK визначає структуру знань}, необхідних для ефективної інтеграції ШІ-інструментів у дизайн-освіту. Технологічний компонент (TK) безпосередньо пов'язаний з технологічним компонентом ІКК, а інтегровані форми знань (TPK, TCK, TPACK) визначають, як ШІ-технології мають бути вплетені у педагогічний процес та предметний зміст;

    \item \textbf{Студійна педагогіка забезпечує методичну основу} для формування ІКК через діяльність. Рефлексивна практика Шона обґрунтовує рефлексивний компонент ІКК, а студійний формат (проєкти, критики, ітерації) створює автентичний контекст для розвитку всіх чотирьох компонентів;

    \item \textbf{Компетентнісний підхід та стандарт В2} визначають нормативні вимоги до результатів навчання і забезпечують відповідність пропонованих змін формальним рамкам вищої освіти. Аналіз прогалин стандарту В2 (таблиця~\ref{tab:b2-ikk-mapping}) визначає конкретні напрямки трансформації.
\end{enumerate}

Ця рамка є наскрізною для всього дисертаційного дослідження: у розділі~\ref{sec:2} вона спрямовує аналіз даних (які аспекти ІКК представлені/відсутні у програмах), а у розділі~\ref{sec:3}~--- обґрунтовує модель трансформації (як саме має змінитися кожна дисципліна для забезпечення формування ІКК).

Важливо підкреслити, що інтегративна рамка не є механічним поєднанням трьох теорій. Вона передбачає їх взаємодію: TPACK інформує про структуру знань, які потрібно сформувати; студійна педагогіка визначає, як ці знання формуються в контексті дизайну; компетентнісний підхід забезпечує відповідність формальним вимогам та валідність результатів. Саме ця тріангуляція теоретичних перспектив відрізняє запропонований підхід від попередніх досліджень, де, як показав огляд~\cite{Musiienko2026ScopingReview}, переважна більшість (73,7\%) залишалися атеоретичними.
